\chapter{Démarche d'ingénierie et conduite du projet}

\section{Un projet guidé par l'usage}

Le point de départ n'était pas le choix d'une base de données ou d'un outil
de visualisation. Il fallait d'abord comprendre ce que chaque utilisateur
chercherait à savoir devant la ligne. Un opérateur a besoin d'une réponse
immédiate : l'imprimante peut-elle marquer la prochaine pièce ? Un technicien
de maintenance travaille autrement. Il compare une dérive, un niveau, une
pression et l'historique des défauts. La direction, enfin, ne souhaite pas lire
des dizaines de tags ; elle veut mesurer l'effet des arrêts sur la production.

Cette différence a orienté tout le projet. Les données sont collectées une
seule fois, mais leur présentation varie selon le rôle. La même logique vaut
pour les indicateurs : avant de calculer une disponibilité, il faut définir
quand la machine est réellement attendue. Ce travail de définition a évité de
construire un tableau techniquement correct mais faux du point de vue de la
production.

\section{Organisation en incréments}

Le développement a suivi des incréments courts. Chacun devait produire une
preuve observable, pas seulement un fichier de configuration. Cette approche
était mieux adaptée qu'un cycle entièrement séquentiel, car deux informations
restaient inconnues : le payload exact de Litmus Edge et la liste définitive
des tags de l'imprimante installée.

\begin{table}[H]
\centering
\caption{Incréments du projet et preuves attendues}
\begin{tabularx}{\textwidth}{p{0.13\textwidth}p{0.30\textwidth}YY}
\toprule
Étape & Travail principal & Preuve produite & Risque retiré \\
\midrule
I1 & Courtier, stockage et Grafana & Services sains et source de données joignable &
Incompatibilité de la pile \\
I2 & Simulateur CIJ & Messages MQTT observables et scénarios reproductibles &
Absence de source industrielle \\
I3 & Flux Node-RED & Points visibles dans InfluxDB et rejets explicites &
Payload invalide ou perdu \\
I4 & Tableaux de bord & Trois vues alimentées en direct &
Indicateurs inutilisables \\
I5 & Alertes et notifier & Alerte de test reçue par webhook &
Supervision uniquement passive \\
I6 & Durcissement et validation & Tests automatisés et procédure de raccordement &
Échec tardif sur site \\
\bottomrule
\end{tabularx}
\end{table}

Une étape n'était considérée comme terminée qu'après un essai sur la
chaîne complète. Par exemple, publier un message MQTT ne suffisait pas pour
valider l'ingestion. Il fallait retrouver le point dans InfluxDB, exécuter la
requête Flux correspondante et afficher la valeur dans Grafana. Cette règle a
mis en évidence les erreurs d'intégration beaucoup plus tôt.

\section{Décomposition du travail}

Le projet a été découpé selon le trajet réel de la donnée. Cette organisation
limite les dépendances entre composants et facilite le diagnostic lorsqu'une
valeur disparaît.

\begin{enumerate}
  \item \textbf{Comprendre la ligne :} identifier le rôle de l'imprimante, les
        situations d'arrêt et les utilisateurs de la supervision.
  \item \textbf{Définir les contrats :} fixer les topics, les champs obligatoires,
        les unités et les identifiants des messages.
  \item \textbf{Produire des données contrôlées :} simuler les états normaux,
        les périodes sans commande, les dérives et les défauts.
  \item \textbf{Transporter et valider :} utiliser MQTT, isoler les messages
        incorrects et garantir un comportement idempotent.
  \item \textbf{Historiser et calculer :} stocker les séries temporelles puis
        ajouter les classifications et prévisions utiles.
  \item \textbf{Restituer et alerter :} adapter les vues aux utilisateurs et
        transmettre les situations qui exigent une action.
  \item \textbf{Prouver :} automatiser les tests, capturer les interfaces
        alimentées et documenter le raccordement futur.
\end{enumerate}

Le dépôt suit cette même décomposition. Le courtier, le simulateur,
l'ingestion, les tableaux de bord et le notifier disposent chacun de leur
dossier. Le fichier Compose reste le point d'assemblage.

\section{Traçabilité entre besoin et validation}

Les exigences du cahier des charges ne sont pas laissées sous forme de liste
isolée. Chaque exigence possède une réalisation identifiable et au moins une
preuve. Cette matrice sert aussi de plan pour la recette sur site.

\begin{longtable}{p{0.10\textwidth}p{0.34\textwidth}p{0.38\textwidth}}
\caption{Matrice de traçabilité simplifiée}\\
\toprule
Réf. & Réalisation & Validation \\
\midrule
\endfirsthead
\toprule
Réf. & Réalisation & Validation \\
\midrule
\endhead
F01 & Trois entrées MQTT dans Node-RED & Publication réelle sur chaque famille de topics \\
F02 & Classification et panneau d'état opérateur & Changement d'état visible après une publication \\
F03 & Champs de niveaux et courbes temporelles & Requêtes Flux exécutées par l'API Grafana \\
F04 & Lissage du débit et champ \texttt{hours\_to\_empty} & Scénario de baisse puis contrôle de la prévision \\
F05 & Date effective et facteur limitant & Lot ouvert depuis 85 jours avec cinq jours restants \\
F06 & Champ \texttt{line\_demanding} et quatre classifications & Essai avec et sans demande \\
F07 & Mesure \texttt{printer\_event} et Pareto & Injection de défauts avec codes et durées \\
F08 & Mesure \texttt{marking\_event} & DMC, grade et résultat retrouvés dans le bucket \\
F09 & Sept règles provisionnées et webhook & Exécution des chaînes de requête d'alerte \\
F10 & Vues opérateur, maintenance et direction & Chargement authentifié et captures des trois vues \\
\bottomrule
\end{longtable}

\section{Gestion des risques}

Les risques ont été traités selon leur effet sur le projet, pas seulement
selon leur probabilité. L'absence d'accès au site était certaine pendant le
développement ; le bon choix n'était donc pas d'attendre, mais de rendre le
prototype indépendant de cette attente.

\begin{table}[H]
\centering
\caption{Principaux risques et mesures prises}
\begin{tabularx}{\textwidth}{p{0.29\textwidth}p{0.12\textwidth}p{0.12\textwidth}Y}
\toprule
Risque & Prob. & Impact & Traitement \\
\midrule
Pas d'accès à Litmus Edge & Forte & Fort & Simulateur et contrat d'entrée stable \\
Tags ou unités différents sur site & Forte & Fort & Couche de mapping concentrée dans Node-RED \\
Faux arrêt pendant une période sans commande & Forte & Fort & Signal de demande obligatoire \\
Doublon provoqué par MQTT QoS~1 & Moyenne & Moyen & Clé InfluxDB stable et essai de retransmission \\
Rupture de communication & Moyenne & Fort & Last Will, détection d'absence et alerte \\
Fuite d'un secret dans le dépôt & Faible & Fort & Variables d'environnement et fichiers exclus \\
Saturation par forte cardinalité & Faible & Moyen & Identifiants volatils stockés comme champs, pas comme tags \\
Alertes trop nombreuses & Moyenne & Fort & Temporisation, regroupement et filtrage par demande \\
\bottomrule
\end{tabularx}
\end{table}

\section{Choix consignés et alternatives écartées}

Plusieurs choix ont changé au cours du travail. Ils ont été conservés dans
les fichiers du projet et dans la documentation afin que la raison reste
compréhensible lors d'une reprise.

\begin{table}[H]
\centering
\caption{Décisions d'architecture structurantes}
\begin{tabularx}{\textwidth}{p{0.25\textwidth}p{0.28\textwidth}Y}
\toprule
Décision & Alternative examinée & Motif \\
\midrule
InfluxDB comme stockage temporel & Base relationnelle avec extension temporelle &
Modèle naturel pour les mesures, écriture Line Protocol et requêtes Flux \\
Node-RED pour l'ingestion & Service Python ou Telegraf &
Mapping visible, modifiable sur site et supervision intégrée \\
Mosquitto sur un nœud & Courtier distribué &
Volume faible et absence de besoin de haute disponibilité au prototype \\
JSON avant Sparkplug B & Protobuf dès la simulation &
Diagnostic plus simple tant que le payload Litmus réel n'est pas connu \\
Trois tableaux de bord & Une vue universelle &
Questions et niveaux de détail différents selon le rôle \\
Notifier derrière un webhook & Canal codé dans chaque alerte &
Changer de canal sans modifier les règles Grafana \\
\bottomrule
\end{tabularx}
\end{table}

\section{Critère de fin d'une fonctionnalité}

Pour ce projet, une fonctionnalité est terminée lorsque cinq conditions sont
réunies : sa configuration est versionnée, son entrée est contrôlée, son
résultat est visible, un test reproductible couvre le risque principal et la
documentation explique son exploitation. Cette définition évite de confondre
une démonstration ponctuelle avec un composant transmissible à une autre
personne.

La commande \texttt{make verify} matérialise ce critère. Elle regroupe les
contrôles de syntaxe, les tests du simulateur et du notifier, le chargement des
tableaux de bord, l'exécution des requêtes Flux et la présence de données ayant
traversé Node-RED. Le chapitre de validation détaille ces preuves.
